  
\addcontentsline{toc}{chapter}{\twkkai{对话录·其七}}  
\markboth{\twkkai{对话录·其七}}{}  
\pagecolor{black}
\color{white}
\quad\quad \twkkai{\lettrine[,lraise=0.2,lhang=0.5]{\textbf{猫}}{：终于回来了……诶，这狄利克雷，一大早又去哪了？要不我先跳支舞迎接他，嘿嘿！}

（猫踮起脚，悄悄地又跳起了舞\marginnote{}{\small ♪}。）

\rule{0pt}{5pt}

\noindent\hspace*{7.4cm}\includegraphics[height=11mm]{frontmatter/猫.pdf}

（狄利克雷踏进了大门。）

狄利克雷：你又在胡闹什么！刚送你回来，没想到你又在捣乱！

猫：没有啦，这不是奏出最后一首曲子来欢迎你嘛……这可是大名鼎鼎的韦伯Op.73第三乐章呢！

\clearpage
\thisblankpage
\includepdf[
  pages={1-3},
  nup=1x3,
  scale=0.7,
  pagecommand={
    \begin{tikzpicture}[remember picture,overlay]
      \node[align=center]
      at ($(current page.north west)+(1.5cm,-11.7cm)$)
      {
        {\Large ♪}\\[0.8em]
        \qraudio
          {weber3_1}
          {\WeberSound}
          {https://raw.githubusercontent.com/kelvinlamresearchmath-web/finding-dirichlet-site/master/weber3_1.mp3}
      };
    \end{tikzpicture}
  }
]{weber3_fixed_v3.pdf}

狄利克雷：刚刚在人间一起和你见了见守门人。你也看到了，他跟我汇报了些东西；所以我现在没有心情跟你闹。

猫：你说说看嘛，汇报了什么——毕竟我在人间是听不懂人话的。或许他所汇报的就能和韦伯这一曲终章产生联系了呢！

狄利克雷：他说他在人间发现了一套理论——已经可以用所谓“类域论”来重新阐释我发现的算术级数定理了……

猫：喵……你好像不兴奋？

狄利克雷：当然。虽然能看着我的符号孩子们长大，我挺高兴；但我隐隐约约感到被威胁了。

猫：哈哈，倒不能这么说——数学总是向前发展的嘛。

守门人：……设 \(K / F\) 是伽罗瓦扩张，且 \(P_{F, \mathfrak{m}}^{+}<H<I_{F}(\mathfrak{m})\)。假设存在某个素理想集合 \(T \subseteq H\)，使得 \(\mathcal{S}_{K / F} \approx T\)。那么
\[
\left[\mathcal{I}_{F}(\mathfrak{m}): \mathcal{H}\right] \leq[K: F],
\]
并且当 \(\chi \neq \chi_{0}\) 且 \(\chi\) 在 \(H\) 上为平凡时，\(L_{\mathfrak{m}}(1, \chi) \neq 0\)……

猫：哎哟，这不是……？

狄利克雷：还要气我对不对？

猫：没有，真的没有！我在想一些事情呢。

狄利克雷：你最好说的是类比。

猫：是的，就是类比！你也知道，在这里，我们的话题其实也仅限于这些了。实际上，这一乐章的速度标记是\textit{Rondo}呢——你可知道“Rondo”之意吗？

狄利克雷：不知道，也不想知道。至少不是现在；不过你想说就说。

猫：好，那我尽量谨慎点。发现了吗——这首简直就是典型回旋曲式，欢快轻盈的主要主题与一个或多个对比主题交替出现着，呼应着；而重点是——这种曲式的核心特点是：一个主题（叠部，refrain）会反复出现，但每一次回归都当然不会是原样重复——它会被装饰、变形、放进不同的和声或调性语境里，周围环绕着不断变化的插部（episode）。听众认得出这个主题，但每次听到时，它的“模样”都不太一样喵。

狄利克雷：又说得头头是道。你是不是在暗示什么？

猫：暗示什么？

狄利克雷：你不觉得你在暗语伤人吗？你说的简直就是我的孩子回旋回来超越我、覆盖我的场景！

猫：不敢，不敢！求求你息怒！

狄利克雷：好。你继续跳吧——我倒是有点喜欢这个乐章了；从头开始吧。

（猫听话、谨慎地迈开了脚步，终于舞起了天庭的最后一曲……）
\clearpage
\begin{center}
\lily[hpadding=0pt]{
    \override Staff.StaffSymbol.color       = #white
    \override Staff.Clef.color              = #white
    \override Staff.TimeSignature.color     = #white
    \override Staff.KeySignature.color      = #white
    \override Staff.BarLine.color           = #white
    \override Staff.LedgerLineSpanner.color = #white
    \override NoteHead.color                = #white
    \override Stem.color                    = #white
    \override Beam.color                    = #white
    \override Flag.color                    = #white
    \override Rest.color                    = #white
    \override Accidental.color              = #white
    \override Dots.color                    = #white
    \override Slur.color                    = #white
    \override Tie.color                     = #white
    \override Script.color                  = #white
    \override DynamicText.color             = #white
    \override Hairpin.color                 = #white
    \override TextScript.color              = #white
    \override TupletBracket.color           = #white
    \override TupletNumber.color            = #white
  \relative bes'' {
  \clef "treble"
  \time 2/4
  \key g \major
  \transposition bes
  g,,16 -.  a16 -. b16 -. c16 -. d16 -. e16 -. fis16 -. g16 -.
    a16 -. b16 -. c16 -. d16 -. e16 -. fis16 -. g16 -. a16 -.
    b8 -.  b8 -. \acciaccatura { b8 ( } a16 ) g16 a16 b16 g8 -.
    b8 -. g8 r8 \bar "||"}
}
\end{center}


\rule{0pt}{20pt}

\noindent\adjustbox{center}{\includegraphics[scale=1]{trebleclef.preview.pdf}}


}


